\documentclass[10pt,a4paper]{amsart}
\usepackage[margin=19mm]{geometry}
\usepackage{amsmath,amssymb,mathtools}
\usepackage[scr=rsfs,cal=euler]{mathalfa}
\usepackage{xfrac}
\usepackage[unicode,hidelinks,bookmarks=false]{hyperref}
\newcommand{\ie}{i.e.\ }
\newcommand{\cf}{cf.\ }
\newcommand{\resp}{respectively\ }
\newcommand{\Subsection}{\subsection}
\providecommand{\nobreakdash}{\nobreak\hbox{-}}
\setlength{\emergencystretch}{2em}
\multlinegap0pt
\usepackage{amssymb}
\usepackage[shortlabels]{enumitem}
\usepackage{tikz}
\newtheorem{thmx}{Theorem}
\newtheorem{prop}{Proposition}
\newtheorem{notation}{Notation}
\newtheorem{lemma}{Lemma}
\newtheorem{thm}{Theorem}
\newcommand{\BD}{{\sf{BD}}}
\newcommand{\Nb}{\mathbb{N}}
\newcommand{\Qb}{\mathbb{Q}}
\newcommand{\Zb}{\mathbb{Z}}
\title{Nonnuclear Bunce--Deddens algebras}
\author{Jamie Bell}
\date{Source: arXiv:2607.28597v1 (CC BY 4.0)}
\begin{document}
\maketitle

\noindent\emph{Source and licence.} Nonnuclear Bunce--Deddens algebras. Version arXiv:2607.28597v1 (CC BY 4.0), distributed by arXiv under the Creative Commons Attribution 4.0 International licence, http://creativecommons.org/licenses/by/4.0/, as recorded in the arXiv licence metadata for this exact version and shown on the abstract page https://arxiv.org/abs/2607.28597v1. Full licence notice: \texttt{licenses/arxiv-2607.28597-CC-BY-4.0.txt}.\par\smallskip

\noindent\textit{Editorial scope.} Counted unit: the article's thm T:Kthy (source0.tex lines 560--569). The article's complete original proof of it is delivered with it (source0.tex lines 571--581, one \texttt{proof} environment of 11 lines). No mathematics is written, repaired or re-derived: the statement and the proof are the article's own lines, in the article's own notation, and no retained line is substituted or re-flowed.  Retained uncounted as the source-local context the proof needs: lines 104--113; lines 149--156; lines 426--435; lines 541--546; lines 548--554.  Nothing was omitted from any retained span, so the package carries no omission record.  No retained line contains a citation, so the wrapper carries no bibliography and no bibliography entry.  No retained line contains a figure environment, an image-inclusion command or drawing code, so no image is delivered and the package declares no asset dependency.  Editorial formatting only, in addition: an AMS \texttt{amsart} preamble that restates the article's own declarations and macros used by the retained lines, namely the environments \texttt{thmx}, \texttt{prop}, \texttt{notation}, \texttt{lemma}, \texttt{thm} on separate counters, together with the standard packages those lines need.  The retained environments print in this document as Thmx 1, Proposition 1, Notation 1, Lemma 1, Theorem 1 in order of appearance, whereas the article prints the same items with the numbering of its own full document.  The complete native source of the article as distributed in the arXiv e-print for this exact version is bundled unchanged as \texttt{sources/206428/source0.tex}.
\begin{thmx}\label{T:B}
    Fix $d\ge 2$. Let $\sigma = (G_n)_{n\in\Nb}$ be a separating normal chain in $F_d$. Then, 
    \[
        K_0(\BD(F_d,\sigma)) \cong \bigcup_{n\in\Nb} \frac{1}{[F_d : G_n]}\Zb\subseteq \Qb,
    \]
    via an isomorphism sending $[1]_0$ to $1$, and 
    \[
        K_1(\BD(F_d,\sigma)) \cong \varinjlim \Zb^{1 + [F_d:G_n](d-1)}. 
    \]
\end{thmx}

\begin{prop}\label{P:equi}
    Let $G\curvearrowright X$ be an action of a countable discrete group on a metrisable Stone space. The following are equivalent:
    \begin{enumerate}
        \item the action is equicontinuous;
        \item there exists a compatible metric on $X$ with respect to which the action is isometric;
        \item there exist finite sets $(X_n)_{n\in \Nb}$, actions $G\curvearrowright X_n$, and $G$-equivariant connecting maps $\varphi_n : X_{n+1} \to X_n$ such that $G\curvearrowright X$ is conjugate to $G\curvearrowright \varprojlim X_n$.
    \end{enumerate}
\end{prop}

\begin{prop}\label{P:limit}
    Let $G\curvearrowright X$ be an equicontinuous action of a countable discrete group on a metrisable Stone space. Then identifying $G\curvearrowright X$ with $G\curvearrowright \varprojlim X_n$ according to Proposition~\ref{P:equi}, we have  
    \[
        C(X)\rtimes_\lambda G \cong \varinjlim \bigoplus_{[y]\in X_n/G} M_{|Gy|}(C^*_\lambda(G_y)).
    \]
    In particular, if $\sigma = (G_n)_{n\in\Nb}$ is a separating normal chain in $G$, then 
    \[
        \BD (G,\sigma) \cong \varinjlim M_{[G:G_n]}(C^*_\lambda(G_n)).
    \]    
\end{prop}

\begin{notation}\label{N:chain}
For a separating normal chain $\sigma = (G_n)_{n\in \Nb}$ in $F_d$, put $m_n = [F_d: G_n]$ and $q_n = [G_n : G_{n+1}] = m_{n+1} / m_n$. Define 
\[
    \Qb_\sigma = \bigcup_{n\in \Nb} \frac{1}{m_n}\Zb \subseteq \Qb. 
\]
\end{notation}

\begin{lemma}\label{L:fin_stage}
    Let $G_n \le F_d$ be a subgroup, and write $m_n := [F_d:G_n]$. Then  
    \[
        K_0(C(F_d/G_n)\rtimes_\lambda F_d) \cong \Zb \quad \text{and} \quad K_1(C(F_d/G_n)\rtimes_\lambda F_d) \cong \Zb^{1+m_n(d-1)}.  
    \]
    Under the identification $K_0(C(F_d/G_n)\rtimes_\lambda F_d) \cong K_0(C^*_\lambda(G_n)) \cong \Zb$, the unit in $C(F_d/G_n)\rtimes_\lambda F_d$ corresponds to $m_n \in \Zb$.
\end{lemma}

\begin{thm}[Theorem~\ref{T:B}]\label{T:Kthy}
    Fix $d\ge 2$. Let $\sigma = (G_n)_{n\in\Nb}$ be a separating normal chain in $F_d$. Then, using Notation~\ref{N:chain}, 
    \[
        K_0(\BD(F_d,\sigma)) \cong \Qb_\sigma \subseteq \Qb, 
    \]
    via an isomorphism sending $[1]_0$ to $1$, and, writing $\beta_n$ for the map in $K$-theory induced by $C(F_d/G_n)\rtimes_\lambda F_d \to C(F_d/G_{n+1})\rtimes_\lambda F_d$,
    \[
        K_1(\BD(F_d,\sigma)) \cong \varinjlim (\Zb^{1 + m_n(d-1)},\beta_n).
    \] 
\end{thm}

\begin{proof}
    Applying Proposition~\ref{P:limit} to decompose $\BD(F_d,\sigma)$ into an inductive limit, by continuity of $K$-theory, we have, for $i\in \{0,1\}$, 
    \begin{equation}\label{E:cty}
        K_i(\BD(F_d,\sigma)) \cong \varinjlim K_i(C(F_d/G_n)\rtimes_\lambda F_d).
    \end{equation}
    Lemma~\ref{L:fin_stage} gives $K_0(C(F_d/G_n)\rtimes_\lambda F_d) \cong \Zb$. Since the connecting maps are unital, they are uniquely determined by the image of the unit. At the $n$th stage, the unit corresponds to $m_n \in \Zb$ and is sent to $m_{n+1} \in \Zb$. It follows that the connecting map $K_0(C(F_d/G_{n})\rtimes_\lambda F_d) \to K_0(C(F_d/G_{n+1})\rtimes_\lambda F_d)$ is given by multiplication by $q_n$. Therefore, 
    \[
        K_0(\BD(F_d,\sigma)) \cong \varinjlim (\Zb \stackrel{\times q_1}{\longrightarrow} \Zb \stackrel{\times q_2}{\longrightarrow} \Zb \longrightarrow \cdots) \cong \Qb_\sigma,
    \]
    where the copy of $\Zb$ at stage $n$ is identified with $\frac{1}{m_n}\Zb \subseteq \Qb$. Since the unit $[1]_0$ corresponds at each stage with $m_n \in \Zb$, it is mapped under this identification to $1 \in \Qb$. The $K_1$-group computation follows from Lemma~\ref{L:fin_stage} and (\ref{E:cty}). 
\end{proof}

\end{document}
